% !PRINTED 276
\weberpage{276}
This group is real when $n$ is a divisor of $1/2\,(p-1)$, because
then and only then the exponents of $\gamma$ are divisible by $p+1$.
The group is contained in $\Gamma_p$ and at the same time in $\Lambda_p$ (\S.\,69)
when $\gamma^{\frac{p^2-1}{2n}}$, $\gamma^{-\frac{p^2-1}{2n}}$
are conjugate, hence when
\[
\gamma^{-\frac{p^2-1}{2n}} = \gamma^{p\frac{p^2-1}{2n}}
\quad \text{or} \quad
\gamma^{(p+1)\frac{p^2-1}{2n}} = 1
\]
that is, when $n$ is a divisor of $1/2\,(p+1)$. This agrees with
\S.\,68, according to which no cyclic groups occur other than those
whose degree is equal to $p$ or a divisor of $1/2\,(p-1)$ or $1/2\,(p+1)$.
In the remaining polyhedral groups $h = 3$, and
we have:

II. \gesperrt{Dihedral groups} of degree $2m$,\quad
$\begin{aligned} \nu &= \nu' = 2,\ \nu'' = m \\ \mu &= \mu' = m,\ \mu'' = 2 \end{aligned}$
\[
D_m = \begin{pmatrix} \gamma^{\frac{p^2-1}{2m}\lambda}, & 0 \\[2pt] 0, & \gamma^{-\frac{p^2-1}{2m}\lambda} \end{pmatrix},
\begin{pmatrix} 0, & \gamma^{\frac{p^2-1}{2m}\lambda} \\[2pt] -\gamma^{-\frac{p^2-1}{2m}\lambda}, & 0 \end{pmatrix},
\quad \lambda = 0, 1, \ldots, m-1.
\]
This group is contained in $L_p$ when $p \equiv 1 \pmod{2m}$, and
in $\Gamma_p$ when $p \equiv -1 \pmod{2m}$, as in Case I.

Here the second representation of the dihedral group given in \S.\,55
has been chosen, in which the distinction of the cases becomes somewhat
simpler than with the first.

III. \gesperrt{Tetrahedral group,}\quad
$\begin{aligned} \nu &= 2,\ \nu' = 3,\ \nu'' = 3,\ n = 12 \\ \mu &= 6,\ \mu' = 4,\ \mu'' = 4. \end{aligned}$

To represent this group, we observe that in the field
$\mathfrak{E}_{2,p}$ an $8$th root of unity
$\gamma^{\frac{p^2-1}{8}}$ always exists. We thus obtain, by \S.\,56, the
three generating substitutions
\[
\Theta = \begin{pmatrix} \gamma^{\frac{p^2-1}{4}}, & 0 \\[2pt] 0, & \gamma^{-\frac{p^2-1}{4}} \end{pmatrix},
\quad
\psi = \begin{pmatrix} 0, & \gamma^{\frac{p^2-1}{4}} \\[2pt] -\gamma^{-\frac{p^2-1}{4}}, & 0 \end{pmatrix},
\quad
\psi_1 = \psi\,\Theta = \begin{pmatrix} 0, & 1 \\ -1, & 0 \end{pmatrix}
\]
\[
\chi = \begin{pmatrix} \dfrac{1}{\sqrt{-2}}\,\gamma^{\frac{p^2-1}{8}}, & \dfrac{1}{\sqrt{-2}}\,\gamma^{\frac{p^2-1}{8}} \\[8pt] \dfrac{1}{\sqrt{-2}}\,\gamma^{-\frac{p^2-1}{8}}, & -\dfrac{1}{\sqrt{-2}}\,\gamma^{-\frac{p^2-1}{8}} \end{pmatrix}.
\]
1) If $p \equiv 1 \pmod 4$, then $\gamma^{\frac{p^2-1}{4}}$ is real,
and $\gamma^{\frac{p^2-1}{8}}$ is real or purely imaginary according
as $p \equiv 1$ or $\equiv 5 \pmod 8$.
